\cotefile{156-6}
\scanpage{50}
\mspage{49}

\uline{Lemme} Soient $I,I'$ deux ens \deleted{\unsure{part.}}\deleted{\unsure{pré}} ordonnés, Alors l'application
\[
\operatorname{Hom}_{\unsure{\mathrm{croiss}|\mathrm{Croiss}}}(I,I')\longrightarrow
\operatorname{Hom}_{\mathrm{cr}}(\unsure{Pf|P_f}(I'),Pf(I))
\]
entre applications croissantes de $I$ dans $I'$, et appl. cr. de $Pf(I)$ en $Pf(I')$ (parties \deleted{\unsure{ouv.}} fermées de $I,I'$) est \uline{injective}. Son image est formée \deleted{des} de certaines des
\begingroup\csname tagsleft@true\endcsname
\[
Pf(I')\longrightarrow Pf(I)\tag*{$(\unsure{x|\alpha})$}
\]
\endgroup
\unsure{qui}, commutent aux intersections \inserted{et réunions} quelconques \deleted{et aux réunions finies \unsure{(la dernière est kif-kif)}}\unsure{,} et \inserted{\unsure{de} toutes si $I$ est sobre (\uuline{p.ex. $I$ fini})}
% slanted
\margin{Elles commutent alors aux réunions quelconques, et inversement}

\deleted{\unsure{Bien}} Passant aux complémentaires ouverts, \inserted{on voit que} les applications $(x)$ correspondent aux \inserted{\unsure{classes d'is.} des} morphismes \deleted{qui \unsure{sont compatibles aux} \unsure{produits|précédents}} de Topos $\operatorname{Top}(I)\unsure{-|\to}\operatorname{Top}(I')$, \deleted{Mais \unsure{ils ne} \illegible{1}} \inserted{ayant une} propriété suppl. (compatibilité aux intersections \uline{quelconques}) \deleted{$I$ est \unsure{évidemment} \unsure{i.e.} sobre,} Ces \deleted{\unsure{appl}} alors correspondent aux applications continues i.e. croissantes
\[
I\longrightarrow\widehat{I}'
\]
où $\widehat{I}'$ est l'espace sobre associé à $I'$, \unsure{i|;}

\uline{NB} on n'avait pas \unsure{$\prime\prime$|à} supposer $I$ séparé ) \deleted{\unsure{on sait} \unsure{$I$}} \deleted{Comme $\unsure{I'^{\circ}|I'^{0}}\longrightarrow\widetilde{I}'$ ($I'$ séparé) cela donne l'injectivité, et pour la bijectivité \unsure{alors|aussi} quand $I'$ sobre i.e. $I'=\widetilde{I}'$}
% slanted
\margin{transformant les $i\in I$ en des \deleted{\unsure{éléments}} \uline{\unsure{stricts|strict}} de $I'$ \inserted{parties fermées irréductibles donc de la forme $\overline{\{i'\}}=I'_{\leq i'}$ (comme $I'$ séparé, cela signifie $I-I'$ \unsure{quel|que|quelc.}}}

Donc cela nous fournit bien (1.55) \deleted{\unsure{pour}} \deleted{au moins si $I'$ sobre \unsure{p.ex.} $I'$ \unsure{fini}} Mais aura-t-on réciproquement que la top de $\widetilde F$ est image inv. de celle de $\widetilde F'$ ? En fait, c'est \uline{équivalent} à la surjectivité de (1.54).

\uline{Ex} Prenons des ateliers ``triviaux'' (p. 44) définis par $(M,\leq,\ll\overset{\mathrm{dfn}}{=}\leq,F)$. \deleted{\unsure{autrement}}\inserted{\unsure{Simplement|Signalons} ici} \deleted{qui \unsure{fait}) \unsure{$\mathcal F$}} : l'axiome At 8, il signifie que pour tt $X\in M$, $\exists\,x\in M$, $x$ \uline{minimal}, tel que $x\leq X$, \unsure{mais} et que \inserted{dans le cas $X\ne x$ i.e. $X$ non \unsure{mini|minim.}} $X$ soit ``successeur'' de
